% page 566 of the facsimile (image p-565 of the scan)
% block 154.9 x 237.7 mm ; left margin 8.0 mm
\pgc{-12.700mm}{0.000mm}{
\l{}
\l{\cel{2}558.}
\l{}
\l{\sou{sturmar}. (\sou{netrans}.) Verbo qua indikas ke atmosfero pertur-}
\l{\cel{3}besas da vento violentoza, pluvo o grelo, fulmini e}
\l{\cel{3}tondri. - DE.}
\l{}
\l{\sou{sturno}. (\sou{zool.}) Pasero korno-rostra,kun plumaro nigra e}
\l{\cel{3}blanka, qua vivas bande en la boski ed en la gardeni.}
\l{\cel{3}- L.\cel{1}\sou{sturnus vulgaris}.\cel{2}FIS.}
\l{}
\l{\sou{stuvo}. Chambro, chamgrego, en qua la temperaturo di la}
\l{\cel{3}aero esas tre alta (multa-kaze : por desinfekto di la}
\l{\cel{3}enklozaji). - EFIS.}
\l{}
\l{su.\cel{2}(\sou{gram.}) Pronomo reflektiva, qua referas la subjekto}
\l{\cel{3}di la propoziciono di qua lu esas un ek la vorti, kande}
\l{\cel{3}la subjekto esas \textquotedbl{}lu\textquotedbl{} o to quon \textquotedbl{}lu\textquotedbl{} reprezentas.- DFIS.}
\l{}
\l{\sou{sub}.\cel{2}(\sou{gram.}) Prepoziciono : l. Indikas la situeso di kozo}
\l{\cel{3}relate to quo esas plu adsupre,\cel{1}\sou{kontaktante}\cel{1}kun lu, en}
\l{\cel{3}la sama direciono vertikala, (a) relate to quon lu su-}
\l{\cel{3}portas, (b) relate to quo kovras lu tote o parte. (\sou{anke}}
\l{\cel{3}\sou{metaf.}) II. Indikas la situeso di kozo relate to quo}
\l{\cel{3}esas plu adsupre,\cel{1}\sou{sen kontaktar}\cel{1}lu, en la sama direcio-}
\l{\cel{3}no vertikala. (\sou{anke metafore)}. - DEFI.}
\l{}
\l{\sou{subalterna}. Inferiora pri lua situeso en hierarkio. - DEFIS}
\l{}
\l{\sou{subisar}.\cel{2}(\sou{trans.}) Tolerar malgre su. FI.}
\l{}
\l{\sou{subita}. Qua eventas neexpektite. - FIS.}
\l{}
\l{\sou{subjekto}. I. (\sou{gram.})e (\sou{logiko}). Termino di la propoziciono}
\l{\cel{3}di qua onu afirmas eso-maniero. - DEFIRS.}
\l{}
\l{\sou{subjuntivo}. (\sou{gram}.) (en\cel{1}\sou{ula lingui nacionala}). Modo di la}
\l{\cel{3}konjugo-sistemo, qua reprezentas la stando o la ago ex-}
\l{\cel{3}presata da verbo qua dependas de altra verbo. DEFIS.}
\l{}
\l{\sou{sublima}. I.Qua esas situita tre alte. - II. (\sou{metaf.}) Qua}
\l{\cel{3}expresas lo bela segun lua formo maxim alta, partikulare}
\l{\cel{3}en la domeni pensala, etikala o morala. - DEFIS.}
\l{}
\l{\sou{submersar}. (\sou{trans.})\cel{2}Plunjar komplete adsub la surfaco di}
\l{\cel{3}aquo (di maro, di lago, di fluvio, di recipiento, ec.)}
\l{\cel{3}EFIS.}
\l{}
\l{\sou{submisar.}\cel{1}(\sou{trans.}) Dependigar de la autoritato di ulu.}
\l{\cel{3}EFIS.}
\l{}
\l{subordinar.- (trans., ad). Dependigar de to quo esas en rango}
\l{\cel{3}plu alta, superiora. - (gram.) Propoziciono subordinita :}
\l{\cel{3}qua dependas, en la sintaxo, de altra propoziciono. - II.}
\l{\cel{4}Dependigar de la ageso o faceso di ulo. - DEFIRS.}
}
